% ─── Capítulo 1: Antecedentes ─────────────────────────────────────────────────
\chapter{ANTECEDENTES DE LA INVESTIGACIÓN}

\section{Antecedentes}

El desarrollo de compiladores ha evolucionado desde procesos manuales de traducción hasta sistemas automatizados capaces de analizar, optimizar y generar código de manera eficiente. Dentro de este proceso, las fases de análisis léxico y sintáctico cumplen un papel esencial como punto de partida, pero un compilador completo requiere además análisis semántico, generación de código intermedio y optimización para garantizar que el programa no solo sea gramaticalmente correcto, sino también coherente y procesable.

Herramientas como Lex/Flex y Yacc/Bison han sido ampliamente utilizadas en el ámbito académico y profesional para la construcción de analizadores. Flex facilita la definición de patrones mediante expresiones regulares, mientras que Bison permite especificar gramáticas y reglas de producción. Sobre esta base, las fases posteriores del compilador permiten validar reglas semánticas, construir representaciones intermedias y aplicar optimizaciones antes de producir el resultado final.

\section{Definición del Lenguaje TPP}

\subsection{Objetivos del Lenguaje}

El lenguaje \textbf{TPP} (\textit{Tiny Programming Processor}) fue diseñado con los siguientes objetivos pedagógicos:

\begin{itemize}
    \item Permitir que los estudiantes de sistemas escriban programas imperativos en un lenguaje cuya sintaxis está en \textbf{español}, reduciendo la carga cognitiva del idioma extranjero al aprender teoría de compilación.
    \item Ser lo suficientemente completo para ilustrar todas las fases del pipeline de compilación (léxico, sintáctico, semántico, código intermedio, optimización y generación de código).
    \item Generar código ejecutable en hardware real (simulado en \textbf{Logisim Evolution}) a través de la ISA RISC-V RV32I.
\end{itemize}

\subsection{Características Sintácticas}

La sintaxis de TPP sigue un estilo imperativo similar a C/Java pero con palabras reservadas en español. A continuación se presenta el resumen de sus construcciones:

\begin{table}[H]
\centering
\begin{tabular}{llp{6cm}}
\toprule
\textbf{Categoría} & \textbf{Construcción} & \textbf{Ejemplo en TPP} \\
\midrule
Tipos de datos & Entero, Decimal & \texttt{entero x = 5;} \\
Arreglos & Estáticos de tamaño fijo & \texttt{entero arr[10];} \\
Control de flujo & si/sino & \texttt{si (x > 0) \{ ... \} sino \{ ... \}} \\
Bucles & mientras, para & \texttt{mientras (i < 10) \{ i = i + 1; \}} \\
Selección múltiple & depende/caso/otros & \texttt{depende (x) \{ caso 1: \{ ... \} \}} \\
Funciones & Con parámetros y retorno & \texttt{fun suma(entero a, entero b) -> entero \{ ... \}} \\
Entrada/Salida & imprimir, leer & \texttt{imprimir("Hola mundo");} \\
Comentarios & Línea simple & \texttt{// Este es un comentario} \\
Operadores arit. & +, -, *, /, \% & \texttt{x = a * 2 + b / 3;} \\
Operadores relac. & ==, !=, <, >, <=, >= & \texttt{si (a != b) \{ ... \}} \\
Retorno de función & retornar & \texttt{retornar a + b;} \\
\bottomrule
\end{tabular}
\caption{Sintaxis del lenguaje TPP}
\label{tab:sintaxis-tpp}
\end{table}

\subsection{Restricciones del Lenguaje}

\begin{itemize}
    \item Los identificadores solo pueden contener letras y dígitos (sin guión bajo \_). Se recomienda \textit{camelCase}.
    \item No existe el tipo booleano explícito; las comparaciones devuelven \texttt{0} (falso) o \texttt{1} (verdadero).
    \item Los arreglos deben declararse con tamaño constante conocido en tiempo de compilación.
    \item No hay recursión completa (modelo de memoria estático sin frames dinámicos).
\end{itemize}

\section{Especificación Léxica}

\subsection{Palabras Reservadas}

Las palabras reservadas de TPP son las siguientes y no pueden usarse como identificadores:

\begin{center}
\texttt{fun \quad si \quad sino \quad mientras \quad para \quad retornar \quad imprimir \quad leer \quad entero \quad decimal \quad depende \quad caso \quad otros}
\end{center}

\subsection{Clases de Tokens}

\begin{table}[H]
\centering
\begin{tabular}{llll}
\toprule
\textbf{Clase} & \textbf{Expresión Regular} & \textbf{Token} & \textbf{Ejemplo} \\
\midrule
Entero & \texttt{[0-9]+} & \texttt{NUM\_ENTERO} & \texttt{42} \\
Decimal & \texttt{[0-9]+\textbackslash.[0-9]+} & \texttt{NUM\_DECIMAL} & \texttt{3.14} \\
Cadena & \texttt{"([\textasciicircum "\textbackslash\textbackslash]|\textbackslash\textbackslash.)*"} & \texttt{CADENA} & \texttt{"hola"} \\
Identificador & \texttt{[a-zA-Z][a-zA-Z0-9]*} & \texttt{ID} & \texttt{miVar} \\
Suma & \texttt{+} & \texttt{MAS} & \texttt{+} \\
Resta & \texttt{-} & \texttt{MENOS} & \texttt{-} \\
Multiplicación & \texttt{*} & \texttt{MULT} & \texttt{*} \\
División & \texttt{/} & \texttt{DIV} & \texttt{/} \\
Igual & \texttt{==} & \texttt{IGUAL} & \texttt{==} \\
Diferente & \texttt{!=} & \texttt{DIFERENTE} & \texttt{!=} \\
Flecha retorno & \texttt{->} & \texttt{FLECHA} & \texttt{->} \\
\bottomrule
\end{tabular}
\caption{Especificación léxica del lenguaje TPP}
\label{tab:especif-lexica}
\end{table}

\section{Expresiones Regulares y Autómatas Finitos}

Las categorías léxicas de TPP se definen formalmente mediante \textbf{Expresiones Regulares (ER)}. Flex convierte internamente estas ER en un \textbf{Autómata Finito No Determinista (NFA)} utilizando la construcción de Thompson, y luego aplica la construcción de subconjuntos para obtener un \textbf{Autómata Finito Determinista (DFA)} equivalente optimizado.

Por ejemplo, la ER para un número entero es:

\[
\texttt{DIGITO}^{+} \equiv [0\text{-}9][0\text{-}9]^{*}
\]

Y la ER para un identificador es:

\[
\texttt{LETRA} \cdot (\texttt{LETRA} \mid \texttt{DIGITO})^{*} \equiv [a\text{-}zA\text{-}Z][a\text{-}zA\text{-}Z0\text{-}9]^{*}
\]

El DFA generado puede procesar cadenas de entrada en tiempo $O(n)$ proporcional a la longitud de la entrada, sin retroceder nunca (propiedad de los DFA).

\section{Diseño de la Gramática (BNF)}

La gramática del lenguaje TPP es una \textbf{Gramática Libre de Contexto (CFG)} escrita en formato BNF. A continuación se presentan las reglas más representativas:

\begin{lstlisting}[style=cstyle, caption={Extracto de la gramática en BNF — parser.y}, language={}]
programa        ::= lista_elementos
lista_elementos ::= elemento | lista_elementos elemento

elemento        ::= declaracion_fun
                  | declaracion_var
                  | instruccion

declaracion_fun ::= "fun" ID "(" parametros ")" "->" tipo "{" bloque "}"
                  | "fun" ID "(" parametros ")" "{" bloque "}"

declaracion_var ::= tipo lista_ids ";"
                  | tipo ID "[" NUM_ENTERO "]" ";"

instruccion     ::= estructura_si
                  | estructura_mientras
                  | estructura_para
                  | estructura_depende
                  | asignacion ";"
                  | llamada_funcion ";"
                  | "retornar" expresion ";"
                  | "imprimir" "(" lista_impresion ")"  ";"
                  | "leer" "(" ID ")" ";"

expresion       ::= expresion "+" expresion
                  | expresion "-" expresion
                  | expresion "*" expresion
                  | expresion "/" expresion
                  | expresion "==" expresion
                  | expresion "!=" expresion
                  | expresion "<"  expresion
                  | expresion ">"  expresion
                  | NUM_ENTERO | NUM_DECIMAL | ID | CADENA
                  | llamada_funcion
\end{lstlisting}

\subsection{Eliminación de Ambigüedad}

La gramática de expresiones es inherentemente ambigua (no especifica precedencia). Bison resuelve esto con declaraciones de precedencia explícitas en el archivo \texttt{parser.y}:

\begin{lstlisting}[style=cstyle, caption={Declaraciones de precedencia en parser.y}]
%left  IGUAL DIFERENTE MENOR MAYOR MENOR_IGUAL MAYOR_IGUAL  // menor prec.
%left  MAS MENOS
%left  MULT DIV MOD                                          // mayor prec.

// Resolver el "dangling else":
%nonassoc LOWER_THAN_SINO
%nonassoc SINO
\end{lstlisting}

Esto garantiza que \texttt{3 + 4 * 2} se evalúe como \texttt{3 + (4*2) = 11} y que un \texttt{sino} siempre se asocie con el \texttt{si} más cercano.

\section{Rúbrica de Evaluación}

El proyecto fue evaluado bajo la siguiente rúbrica:

\begin{table}[H]
\centering
\small
\begin{tabular}{rp{7.5cm}cccc}
\toprule
\textbf{N°} & \textbf{Criterio} & \textbf{Excelente} & \textbf{Bueno} & \textbf{Regular} & \textbf{Puntaje} \\
\midrule
1  & Definición del lenguaje (objetivos, sintaxis, características) & 1.0 & 0.75 & 0.50 & 1.0 \\
2  & Especificación léxica (tokens, reservadas, operadores)         & 1.5 & 1.0  & 0.5  & 1.5 \\
3  & Expresiones regulares y autómatas finitos                      & 1.5 & 1.0  & 0.5  & 1.5 \\
4  & Implementación del analizador léxico                           & 2.0 & 1.5  & 1.0  & 2.0 \\
5  & Diseño de la gramática (BNF/EBNF) y eliminación de ambigüedad & 2.0 & 1.5  & 1.0  & 2.0 \\
6  & Implementación del analizador sintáctico                       & 2.0 & 1.5  & 1.0  & 2.0 \\
7  & Implementación del analizador semántico y tabla de símbolos    & 2.0 & 1.5  & 1.0  & 2.0 \\
8  & Generación de código intermedio                                & 1.5 & 1.0  & 0.5  & 1.5 \\
9  & Optimización de código                                         & 1.0 & 0.75 & 0.50 & 1.0 \\
10 & Generación de código final (ensamblador RV32I)                 & 2.0 & 1.5  & 1.0  & 2.0 \\
11 & Funcionamiento integral (pruebas y manejo de errores)          & 1.5 & 1.0  & 0.5  & 1.5 \\
12 & Documentación técnica, presentación y defensa del proyecto     & 2.0 & 1.5  & 1.0  & 2.0 \\
\midrule
   & \textbf{TOTAL} & & & & \textbf{20.0} \\
\bottomrule
\end{tabular}
\caption{Rúbrica de evaluación del compilador TPP}
\label{tab:rubrica}
\end{table}

\section{Objetivos}

\subsection{Objetivo General}

Diseñar e implementar la versión final de un compilador para un minilenguaje estructurado en español, integrando análisis léxico, análisis sintáctico, análisis semántico, generación de código intermedio y optimización de código, con salida en ensamblador RISC-V RV32I ejecutable en Logisim Evolution.

\subsection{Objetivos Específicos}

\begin{itemize}
    \item \textbf{Consolidar las fases iniciales:} Mantener e integrar el analizador léxico en Flex y el analizador sintáctico en Bison como base del compilador del minilenguaje.
    \item \textbf{Implementar el análisis semántico:} Verificar reglas de coherencia del programa, como uso de identificadores, compatibilidad de tipos y validez de instrucciones.
    \item \textbf{Generar código intermedio:} Transformar las instrucciones reconocidas en una representación intermedia (TAC) que facilite su posterior procesamiento.
    \item \textbf{Optimizar el código generado:} Aplicar plegado de constantes, simplificación algebraica y eliminación de código muerto.
    \item \textbf{Generar código RV32I:} Traducir el AST optimizado a instrucciones ensamblador RISC-V ejecutables en el procesador Logisim.
    \item \textbf{Integrar con Logisim:} Utilizar \texttt{compile\_to\_logisim.py} para convertir el ensamblador a formato hexadecimal cargable directamente en la RAM del procesador.
    \item \textbf{Validar el compilador final:} Probar el sistema con programas del minilenguaje que evidencien el funcionamiento integrado de todas las fases.
\end{itemize}
